% Standalone appendix document - compiles on its own:
%   latexmk -pdf 01-technical-analysis.tex
\documentclass[a4paper,11pt]{article}
\usepackage[english]{babel}
\usepackage{../../appendix-style}

\title{2.1 Technical Analysis}
\author{SW4PRJ4 -- Group 3}
\date{\today}

\begin{document}
\maketitle

\begin{versionhistory}
  1.0 & 09-09-2026 & Group 3 & Initial version \\
  1.1 & 16-09-2026 & Group 3 & Stack updated to what is actually set up. Added choice of authentication, deployment, database, version control and git workflow \\
  1.2 & 23-09-2026 & Group 3 & Added operations, CI and security. Authentication status and CI description updated \\
\end{versionhistory}

\section{Purpose}
This appendix describes the technologies we use in PrepMeUp and why we chose them. For the bigger decisions we also list the alternatives we looked at and why we did not pick them. The system consists of one backend, two user interfaces and one shared client package. The requirements are given in appendix 1.2.

\section{Technology Stack}

\subsection{Clients (TypeScript monorepo, \texttt{clients/})}
\begin{itemize}
    \item Customer app: Expo (React Native) with TypeScript and Expo Router. Android is the main target (NFR-06). Expo can also build for the web, which makes it easy to try the app in a browser during development.
    \item Admin panel: React, Vite, TypeScript and React Router. Runs in the browser.
    \item Shared client package \texttt{packages/api-client}: generated from the OpenAPI document of the API with NSwag. It is never edited by hand.
    \item The three packages are one npm workspace. \texttt{npm install} is run once in \texttt{clients/}, and \texttt{npm test} and \texttt{npm run typecheck} run across all of them.
\end{itemize}

\subsection{Backend (.NET)}
\begin{itemize}
    \item ASP.NET Core Web API on .NET 10 with controllers and dependency injection.
    \item Serilog for structured logging.
    \item ProblemDetails for central error handling.
    \item OpenAPI document with Swagger UI in development. The document is the contract between the API and the clients.
    \item CORS with a list of allowed origins (admin panel and customer app) read from configuration.
    \item Health check endpoint on \texttt{/health}.
    \item BackgroundService as scheduler (planned): checks expiry dates and plans new deliveries.
    \item All configuration and secrets come from environment variables.
\end{itemize}

\subsection{Data}
\begin{itemize}
    \item SQL Server 2022 in Docker, accessed through EF Core.
    \item Migration strategy: EF Core Migrations, code-first. The model in \texttt{PrepMeUp.Infrastructure} is the source of truth, and every schema change is a generated migration kept in git. Production applies them as a migration bundle, see the deployment pipeline appendix.
    \item The database port is not exposed to the host. Only the \texttt{api} container can reach it.
    \item Data is stored in a Docker volume, so it survives when the containers are recreated.
\end{itemize}

\subsection{Authentication}
\begin{itemize}
    \item Keycloak 26 as an external identity provider, running in the same docker compose as the API.
    \item One realm \texttt{prepmeup} with two roles: \texttt{household} and \texttt{admin}. Customers can register themselves and get the \texttt{household} role automatically.
    \item Two public clients, \texttt{prepmeup-mobile} and \texttt{prepmeup-admin}. Both use the Authorization Code flow with PKCE. Logging in with username and password directly through the client (password grant) is turned off.
    \item The realm is described in \texttt{infra/keycloak/realm-prepmeup.json} and imported when Keycloak starts. That way every group member gets the same setup.
    \item Status: Keycloak runs, and the API validates the JWT. \texttt{GET /items} is behind \texttt{[Authorize]} as the first protected endpoint. Checking roles and household ownership is not implemented yet (section~\ref{sec:ops}).
\end{itemize}

\subsection{Test}
\begin{itemize}
    \item Backend unit tests: NUnit and NSubstitute against the Application layer.
    \item Backend integration tests: \texttt{WebApplicationFactory} against the API.
    \item Client tests: Vitest.
    \item GitHub Actions builds and tests the backend and the clients on every pull request. The pipeline must be green before merge. On \texttt{main} the same run also builds the images and starts the deployment (section~\ref{sec:ops}).
\end{itemize}

\subsection{Deployment and Development Tools}
\begin{itemize}
    \item Local development: \texttt{docker compose up} starts \texttt{api} (port 5001), \texttt{db} and \texttt{keycloak} (port 8082) on the network \texttt{prepmeup-net}. It works the same on Windows and Linux.
    \item Production: \texttt{docker-compose.prod.yml} deployed with Coolify on a group member's own server (section~\ref{sec:deployment}).
    \item GitHub for version control, pull requests and CI.
    \item Kaneo (self-hosted) as task board.
    \item Claude Code and Antigravity as AI assistants. See section~\ref{sec:ai}.
\end{itemize}

\section{Architecture}
\begin{itemize}
    \item Layered architecture with four projects:
    \begin{itemize}
        \item \texttt{PrepMeUp.Api}: presentation
        \item \texttt{PrepMeUp.Application}: use cases, repository interfaces
        \item \texttt{PrepMeUp.Domain}: entities and rules
        \item \texttt{PrepMeUp.Infrastructure}: EF Core, repositories, migrations
    \end{itemize}
    \item Rule: a layer may only use the layer below it. Domain references nothing.
    \item Infrastructure implements the interfaces defined in Application. Api only knows Infrastructure in the composition root (\texttt{Program.cs}).
    \item Cross-cutting concerns: authentication, logging, error handling and configuration.
    \item The scheduler shares the Application layer with the Web API.
\end{itemize}

\section{Choice of Authentication}
\label{sec:auth}
The system needs login for two kinds of users: customers in the app and administrators in the admin panel. At the meeting on 15-09 we agreed that step 2 of the project is to add an external authority, and we listed Firebase, Keycloak, Clerk and OAuth as options. We also looked at the solution from the course, ASP.NET Core Identity with our own JWT.

``OAuth'' is not really an alternative on the same level as the others. OAuth 2.0 is a standard for how an app gets a token from an authorization server, and OpenID Connect (OIDC) adds login on top of it. Keycloak is a full OAuth 2.0 and OIDC server. Firebase and Clerk also give the app a JWT after login, but they are mostly used through their own SDKs. So the real question was which service should handle login, not whether to use OAuth.

\begin{center}
\begin{tabular}{@{}p{3.1cm} p{5.6cm} p{5.6cm}@{}}
\textbf{Option} & \textbf{For} & \textbf{Against} \\
\hline
ASP.NET Core Identity + own JWT & Taught in the course. Everything stays in our own API and database. & We write login, token refresh, password reset and email verification ourselves. It is not an external authority. \\
\hline
Firebase Authentication & Hosted by Google, quick to get started, good SDKs for mobile. & User data lives at Google. Roles need custom claims set with the Admin SDK. Built around the rest of Firebase, which we do not use. \\
\hline
Clerk & Hosted service with ready-made login screens for React and Expo. & User data lives at a third party. Paid plans above the free tier. Hard to move away from later. \\
\hline
Keycloak & Open source and free. Self-hosted in our own compose. Standard OIDC, so the API only needs the normal JWT bearer setup. Roles are built in. & Heavier to run (Java). More to configure. We have to run and update it ourselves. Only mentioned briefly in the course. \\
\hline
\end{tabular}
\end{center}

We chose \textbf{Keycloak}. The main reasons are:
\begin{itemize}
    \item \textbf{We host the rest ourselves.} Kaneo already runs on our own server, and the API and database are deployed there as well (section~\ref{sec:deployment}). Keycloak is just one more container in the same compose, so it fits the setup we have.
    \item \textbf{Personal data stays with us.} Customers give us their name, address and household. With Keycloak that data does not go to Google or another company, which is simpler in terms of GDPR.
    \item \textbf{Standards instead of a vendor SDK.} The apps log in with Authorization Code + PKCE, which is the flow recommended for mobile and browser apps, and the API only checks a normal JWT. If we later want to change provider, the API code stays almost the same.
    \item \textbf{Roles fit our users.} The \texttt{household} and \texttt{admin} roles are set in Keycloak and end up in the token, so the API can use them for authorization.
    \item \textbf{Free.} No limits on the number of users and no credit card.
\end{itemize}

The downside is that Keycloak is more work to set up and runs as a separate Java service that uses more memory than the API itself. We accepted that, because the realm file makes the setup repeatable, and because learning a real identity provider is part of what we want to get out of the project.

\section{Choice of Deployment}
\label{sec:deployment}
At the meeting on 15-09 we also agreed that we need one shared deployment. When there is an administrator who dispatches deliveries, everyone has to work against the same data. It does not work if each group member only runs the system locally.

\begin{center}
\begin{tabular}{@{}p{3.1cm} p{5.6cm} p{5.6cm}@{}}
\textbf{Option} & \textbf{For} & \textbf{Against} \\
\hline
Local only & Nothing to set up. & No shared data, no real administrator, no test on a real phone outside the home network. \\
\hline
Azure & Used in the course. Managed services. & Costs money or student credits that run out. Many services to set up (App Service, SQL, container registry). \\
\hline
Coolify on own server & Free, since the server already exists. Deploys straight from our docker compose. Handles HTTPS certificates. & Depends on one server and one group member. Not in the course. \\
\hline
\end{tabular}
\end{center}

We chose \textbf{Coolify}. The production setup is in \texttt{docker-compose.prod.yml}. It differs from the local compose in a few ways:
\begin{itemize}
    \item No ports are published. Coolify's reverse proxy handles HTTPS and forwards to the containers.
    \item Keycloak runs in production mode with its own PostgreSQL database (\texttt{keycloak-db}). User accounts are kept apart from the application data in SQL Server.
    \item All public URLs and passwords come from environment variables set in Coolify. The compose file refuses to start if one is missing.
    \item All services restart automatically if they stop.
\end{itemize}

\section{Choice of Database}
We use SQL Server for the application data. It is the database used in the backend course, so the course material and exercises apply directly, and EF Core has full support for it. As the course suggests, the integration tests can use SQLite in memory instead, so they do not need a running SQL Server. PostgreSQL was also considered, but we did not see a reason to move away from what the course uses. Keycloak uses its own PostgreSQL database in production, because it is a separate system with its own data.

\section{Version Control and Git Workflow}
\label{sec:git}
The repository is on GitHub. Coolify reads the production compose file from GitHub and pulls the images from GitHub Container Registry. The course uses GitLab CI, but the ideas are the same, so we use GitHub Actions.

We looked at three workflows from the course: trunk-based, GitHub Flow and Git Flow. We chose \textbf{GitHub Flow}.
\begin{itemize}
    \item The course describes GitHub Flow as a good fit for small teams building web applications that do not need several versions at the same time and release often. That matches us: one version, deployed continuously to Coolify.
    \item Git Flow is meant for planned releases and several supported versions. Its \texttt{develop} and \texttt{release} branches would give us more branches to keep track of without any benefit.
    \item Pure trunk-based development, where everyone commits to \texttt{main}, would skip the code review we want before merge.
\end{itemize}

The rules are:
\begin{itemize}
    \item \texttt{main} is the only long-lived branch and can always be deployed.
    \item Work happens on short-lived branches named \texttt{<type>/<slug>}. The number in for example \texttt{feat/pmu-12-household-profile} is the Kaneo task.
    \item Commit messages follow Conventional Commits, for example \texttt{feat(api): add household endpoint}.
    \item Changes go into \texttt{main} through a pull request with green CI and one review. We use squash merge, so each pull request becomes one commit on \texttt{main}.
    \item Releases are git tags on \texttt{main}, for example \texttt{v0.1.0} at the end of a sprint.
\end{itemize}
Branch names and commit messages are checked by git hooks on each computer and again by GitHub Actions on every pull request. The details are in appendix 10.1.

\section{Operations, CI and Security}
\label{sec:ops}
This section describes how the system is built, deployed and protected. The step-by-step deployment chain, the rollback procedure and the rules for migrations are in appendix 10.3. Here we explain the choices and look at what could go wrong.

\subsection{Containers}
Every service runs as a Docker container, both locally and in production.
\begin{itemize}
    \item The API image is built in two stages, as taught in the course. The first stage uses the .NET SDK image to restore and publish. The second stage only contains the ASP.NET runtime and the published files. The SDK and the source code are not in the image that runs in production.
    \item The API container runs as a non-root user (\texttt{USER \$APP\_UID}). If someone finds a hole in the application, they do not get root inside the container.
    \item The Keycloak image is the official Keycloak image with our realm file copied in. The realm configuration is then part of a versioned image instead of something set up by hand in production.
    \item The base images use version tags, never \texttt{latest} (for example \texttt{keycloak:26.3} and SQL Server \texttt{2022-CU27}), so a new major release upstream does not reach production without a commit.
\end{itemize}

\subsection{CI and the Gate Before Deploy}
The workflow \texttt{ci.yml} has four jobs:
\begin{enumerate}
    \item \texttt{backend}: \texttt{dotnet build} and \texttt{dotnet test}.
    \item \texttt{clients}: \texttt{npm ci}, typecheck and Vitest.
    \item \texttt{publish}: only on \texttt{main}, and only if both jobs above are green. Builds the API and Keycloak images and pushes them to GitHub Container Registry, tagged with the commit SHA and \texttt{main}.
    \item \texttt{deploy}: only if \texttt{publish} is green. Calls Coolify's deploy webhook.
\end{enumerate}
A second workflow, \texttt{conventions.yml}, checks branch names and pull request titles.

The course defines continuous integration as frequent integration where each integration is checked by an automated build that includes the tests. We take that one step further: the automated build is also the only way to get something into production. Because \texttt{publish} depends on the test jobs and \texttt{deploy} depends on \texttt{publish}, a red test means that no image is built and no deployment starts. Automatic deployment on push must be turned off in Coolify (appendix 10.3), so there is no second route. The image that runs in production is the same image that was built after the tests passed; Coolify does not build anything itself.

The deploy job only succeeds if Coolify answers exactly \texttt{200} over HTTPS. The first version only failed on \texttt{4xx} and \texttt{5xx}. A redirect from Cloudflare (\texttt{301}) therefore gave a green job while nothing was deployed. A green pipeline must mean that the deployment really started, otherwise the gate gives false confidence. The details are in appendix 10.3.

\subsection{Network and HTTPS}
\begin{itemize}
    \item Only Coolify's reverse proxy is reachable from the internet. It ends HTTPS and forwards the requests over plain HTTP to \texttt{api} and \texttt{keycloak} on the internal Docker network. The course describes this setup: only the proxy needs the certificate, and the proxy limits what is exposed to the internet.
    \item No container publishes a port in production. \texttt{db} and \texttt{keycloak-db} can only be reached from inside the network. The network layout is in appendix 3.1.
    \item The connection from the API to SQL Server is not encrypted (\texttt{Encrypt=False}). We accept this because the traffic never leaves the internal Docker network on one server.
\end{itemize}

\subsection{Authentication and Authorization in the API}
The API checks every access token it receives. The token must be signed by Keycloak's key, have the expected issuer and the audience \texttt{prepmeup-api}, and not be expired. We allow 30 seconds of clock difference. Access tokens live for 15 minutes.

The API fetches Keycloak's discovery document over the internal network, while the tokens carry the public Keycloak address as issuer. The two addresses are different in production. That is why the issuer is set explicitly and not taken from Keycloak's discovery document. Chapter 8 in the report describes how we found this. Fetching the discovery document over plain HTTP is allowed (\texttt{RequireHttpsMetadata=false}) for the same reason as above: the call stays inside the Docker network.

The requirements NFR-09 to NFR-11 describe what the API must refuse:
\begin{itemize}
    \item \textbf{NFR-11, no token gives 401.} Handled by \texttt{[Authorize]} and the JWT check above. In place for \texttt{GET /items}.
    \item \textbf{NFR-10, admin endpoints give 403 to customers.} The \texttt{admin} role is in the token. The role check on the endpoints is not implemented yet, because there are no admin endpoints yet.
    \item \textbf{NFR-09, a household only sees its own data.} A role cannot answer this. The API must compare the household in the request with the user in the token. This check belongs in the use case in the Application layer, so it cannot be forgotten in a single controller. Not implemented yet.
\end{itemize}

CORS only allows the origins of the admin panel and the customer app, read from configuration. There is no wildcard. CORS only limits what a browser allows. It does not protect the API against other clients, which is why authorization must be checked in the API itself.

\subsection{Secrets}
\begin{itemize}
    \item Production secrets are only set as environment variables in Coolify. They are not in the repository, not in the compose file and not in any image. That is why the images on GitHub Container Registry can be public.
    \item \texttt{docker-compose.prod.yml} marks every required secret with \texttt{:?}. If one is missing, the stack does not start. It does not start with an empty password.
    \item GitHub stores only two secrets: the Coolify webhook URL and an API token that can only start a deployment. The token cannot read the environment variables in Coolify, so it does not leak any passwords.
    \item Pushing images uses the token that GitHub Actions creates for each run. The \texttt{publish} job asks for write access to packages and nothing else beyond reading the repository. There is no long-lived registry password.
    \item AI tools cannot read \texttt{.env} files. This is enforced in \texttt{.claude/settings.json}.
\end{itemize}

\subsection{Weaknesses We Know About}
These are known and accepted for now. They are also listed under open decisions.
\begin{itemize}
    \item The API connects to SQL Server as \texttt{sa}. The API only needs to read and write data, so it should have its own database user with fewer rights. The migrations, which change the schema, would keep the stronger user.
    \item If \texttt{Keycloak:Authority} is empty, the API starts without authentication. A protected endpoint then fails with a server error instead of 401. In production this cannot happen, because the compose file requires the value. The API itself should still refuse to start outside development (NFR-07).
    \item The API does not redirect HTTP to HTTPS and does not send an HSTS header itself. The course says an API should only accept HTTPS. In our setup the proxy handles HTTPS, so this depends on how Coolify's proxy is configured. It has not been checked.
    \item The GitHub Actions in the workflows are pinned to a major version (for example \texttt{actions/checkout@v4}), not to a commit SHA. If one of those actions is compromised, our pipeline runs the new code.
\end{itemize}

\section{AI Tools}
\label{sec:ai}
We use AI assistants as a development tool (decided on 09-09). To keep the work understandable for everyone in the group, we set some limits:
\begin{itemize}
    \item The shared rules are written in \texttt{AGENTS.md}. Claude Code and Antigravity both read that file, so they work from the same rules.
    \item Claude Code has four read-only review agents (architecture, backend, test design, security) and a set of skills for recurring tasks, for example planning, TDD and git workflow. The review agents are copied to Antigravity with a small script, so the two tools do not drift apart.
    \item The AI can read the Kaneo board and create and comment on tasks, but not move, assign or delete them. Planning the sprint is the group's job. This is enforced with a deny list in \texttt{.claude/settings.json}.
    \item The AI cannot read \texttt{.env} files or key files.
    \item Test cases are designed and approved by a person before they are implemented, and whoever commits must be able to explain the change.
\end{itemize}

\section{Key Technical Decisions}
\begin{itemize}
    \item One backend serves both user interfaces. The OpenAPI document is the single contract.
    \item Business rules live in Domain only. Clients hold UI state and call the API through the generated client.
    \item EF Core instead of hand-written SQL. Generated SQL can be inspected and replaced where needed.
    \item The database is not exposed outside Docker.
    \item MediatR is not used. It solves a problem the project does not have and is not part of the curriculum.
    \item Vertical slices: each iteration delivers one feature through all four layers, Domain first.
    \item All user-facing text is in English, with no translation layer (NFR-01).
\end{itemize}

\section{Open Decisions}
\begin{itemize}
    \item Data fetching and token storage in the customer app, for example TanStack Query and Expo SecureStore. Not installed yet.
    \item Whether the customer app is also released as a web app.
    \item A separate database user for the API with only read and write rights, instead of \texttt{sa}.
    \item The API should refuse to start without \texttt{Keycloak:Authority} outside development.
    \item Check whether Coolify's proxy redirects HTTP to HTTPS and sends HSTS for the API and Keycloak.
    \item Whether to pin GitHub Actions to commit SHAs.
\end{itemize}

\section*{Sources}
Coolify and Keycloak integration are not covered by the course material; the course lists Keycloak once as an example of an identity provider. GitHub Actions is not covered either -- the course teaches CI with GitLab CI. The principles used in section~\ref{sec:ops} are from the course:
\begin{itemize}
    \item Multi-stage Docker build and publishing images to a registry: \path{context/bad/markdown/01-intro-og-docker/docker.md}, lines 135--156.
    \item CI as an automated build including tests: \path{context/swt/markdown/02.1-continuous-integration/01-Continuous-Integration.md}, lines 43--50.
    \item Reverse proxy, HTTPS only at the proxy, forwarded headers: \path{context/bad/markdown/11-release-security/release-og-deployment.md}, lines 42--85.
    \item Production secrets only on the production host: same file, lines 125--131.
    \item APIs should only accept HTTPS, HSTS: \path{context/bad/markdown/11-release-security/https.md}, lines 89--110.
    \item JWT claims (issuer, expiration, audience) and signature validation: \path{context/bad/markdown/10-auth/json-web-token.md}, lines 49--52 and 76--129.
    \item Keycloak listed as an identity provider: \path{context/bad/markdown/10-auth/authentication-og-authorization.md}, line 808.
\end{itemize}

\end{document}
